% Part: normal-modal-logic
% Chapter: tableaux
% Section: completeness
%READERNOTE{SOL6-A180: संपूर्णतेच्या रचनेत बंद आणि संपूर्ण शाखा वेगळ्या केल्या. गृहीतकांचे समान पूर्वचिन्ह 1 स्पष्ट केले; रिक्त संचाचे प्रकरण वेगळे हाताळले. सांत रचनेचे थांबणे मोडल खोली आणि सांत संततीवरून स्पष्ट केले. असांत संचासाठी प्रत्येक सांत उपसंचाची पूर्ततायोग्यता, सुसंगततेची सांत सिद्धता, लिंडेनबाउमचे पूर्वप्रमेय आणि कॅनॉनिकल सत्यता पूर्वप्रमेय वापरून सर्वसाधारण निष्कर्ष जपला.}

\documentclass[../../../include/open-logic-section]{subfiles}

\begin{document}

\olfileid{nml}{tab}{cpl}

\olsection{\Log{K} साठी संपूर्णता}

\begin{explain}
  !!{tableau}s ची पद्धत संपूर्ण आहे हे
  दाखवण्यासाठी पुढील विधान सिद्ध करावे लागते:
  $\Gamma \Proves !A$ दाखवणारा बंद !!{tableau}
  नसेल, तर $\Gamma \Entails/ !A$; म्हणजे
  प्रतिदृष्टांत प्रतिमान अस्तित्वात असते.
  पण ‘बंद !!{tableau} नाही’ याचा अर्थ
  तो बांधण्याचा प्रत्येक प्रयत्न बंद होण्यात
  अयशस्वी होतो. युक्ती अशी की प्रत्येक प्रयत्न
  अयशस्वी होत असेल, तर एक विशिष्ट
  \emph{पद्धतशीर आणि सर्वसमावेशक} प्रयत्नही
  अयशस्वी होतो. बंद !!{tableau} अस्तित्वात
  असेल, तर या पद्धतशीर प्रयत्नाने तो बंद
  झाला असता. अशा एकाच टॅब्लोमधील खुल्या
  शाखांत प्रतिदृष्टांत प्रतिमानाची व्याख्या
  करण्यासाठी लागणारी सर्व माहिती असते.
  त्यातील एका खुल्या शाखेवरून मिळणाऱ्या
  प्रतिदृष्टांत प्रतिमानाची जगे म्हणजे त्या
  शाखेवर वापरलेली सर्व पूर्वचिन्हे; आणि
  !!a{propositional
    variable}~$p$ हे~$\sigma$ येथे
  तेव्हा आणि तेव्हाच सत्य असते जेव्हा
  $\sFmla{\True}{p}[\sigma]$ हे त्या
  शाखेवर आलेले असते.
\end{explain}

\begin{defn}
  !!a{tableau} मधील शाखेला संपूर्ण
  म्हणतात, जर नियम लावता येईल असे
  $\sFmla{S}{!A}[\sigma]$ हे पूर्वचिन्हयुक्त
  !!{formula} शाखेवर असताना पुढील गोष्टीही
  तिच्यावर असतील:
  \begin{enumerate}
    \item विधानीय, एकाच शाखेवर सूत्रे जोडणाऱ्या
      नियमाच्या बाबतीत त्याच्या निष्कर्षांतील
      पूर्वचिन्हयुक्त !!{formula}s;
    \item विधानीय, शाखा फोडणाऱ्या नियमाच्या
      बाबतीत संबंधित निष्कर्षरूप
      !!{formula}s पैकी एक;
    \item नवे पूर्वचिन्ह लागणाऱ्या मोडल नियमाच्या
      बाबतीत किमान एक शक्य निष्कर्ष;
    \item वापरलेले पूर्वचिन्ह लागणाऱ्या मोडल
      नियमाच्या बाबतीत शाखेवर असलेल्या
      प्रत्येक योग्य पूर्वचिन्हासाठी संबंधित
      निष्कर्ष.
  \end{enumerate}
\end{defn}

\begin{explain}
उदाहरणार्थ, संपूर्ण शाखेवर
$\sFmla{\True}{!B \land !C}[\sigma]$
असेल, तर $\sFmla{\True}{!B}[\sigma]$
आणि $\sFmla{\True}{!C}[\sigma]$
ही दोन्ही असतात. तिच्यावर
$\sFmla{\True}{!B \lor !C}[\sigma]$
असेल, तर $\sFmla{\True}{!B}[\sigma]$
आणि $\sFmla{\True}{!C}[\sigma]$
यांपैकी किमान एक असते.
शाखेवर \iftag{prvBox}
{$\sFmla{\False}{\Box !B}[\sigma]$}
{$\sFmla{\True}{\Diamond !B}[\sigma]$}
असेल, तर काही~$n$ साठी अनुक्रमे
\iftag{prvBox}{$\sFmla{\False}{!B}[\sigma.n]$}
{$\sFmla{\True}{!B}[\sigma.n]$}
हेही असते. आणि शाखेवर
\iftag{prvBox}{$\sFmla{\True}{\Box !B}[\sigma]$}
{$\sFmla{\False}{\Diamond !B}[\sigma]$}
असेल, तर शाखेवर $\sigma.n$ वापरलेले
असेल अशा प्रत्येक~$n$ साठी अनुक्रमे
\iftag{prvBox}{$\sFmla{\True}{!B}[\sigma.n]$}
{$\sFmla{\False}{!B}[\sigma.n]$}
हेही असते.
\end{explain}

\begin{prop}\ollabel{prop:complete-tableau}
  $\Gamma$ हा $1$ या एकाच पूर्वचिन्हाने चिन्हांकित सूत्रांचा
  सांत संच असू द्या. त्यासाठी असा !!a{tableau} असतो की
  त्याची प्रत्येक शाखा बंद किंवा संपूर्ण असते.
\end{prop}

\begin{proof}
  $\Gamma$ साठीच्या !!a{tableau} मधील
  एखादी खुली शाखा घ्या. त्या शाखेवर नियम
  लावता येतील अशी सांत पूर्वचिन्हयुक्त
  !!{formula}s आहेत. ठरलेल्या एखाद्या क्रमाने
  (उदाहरणार्थ, वरून खाली), ज्या
  पूर्वचिन्हयुक्त !!{formula}s साठी
  (1)--(4) या अटी आधीच पूर्ण होत नाहीत,
  त्यांपैकी प्रत्येकाला लागू होणारे नियम
  लावून शाखा विस्तारित करा. काही वेळा शाखा
  फुटतील; मग उरलेल्या पूर्वचिन्हयुक्त
  !!{formula}s साठी प्रत्येक नव्या शाखेच्या
  टोकावर नियम लावा. पूर्वचिन्हयुक्त
  !!{formula}s आणि शाखेवर वापरलेली
  पूर्वचिन्हे यांची संख्या सांत असल्याने
  या प्रक्रियेतून मूळ शाखेपासून पुढे जाणाऱ्या
  (कदाचित अनेक) शाखा मिळतात. त्यांपैकी
  प्रत्येकावर ही प्रक्रिया पुन्हा लावा.
  बंद झालेल्या शाखांवर पुढील नियम लावू नयेत.
  ही प्रक्रिया सांत पायऱ्यांत थांबते: नव्या पूर्वचिन्हावर येणारी
  सूत्रे त्याच्या पालकावरील मोडल सूत्रांची उपसूत्रे असतात.
  त्यामुळे प्रत्येक पूर्वचिन्ह-विस्ताराने उपलब्ध मोडल खोली कमी
  होते; प्रारंभीच्या सांत सूत्रसंचाची कमाल मोडल खोली ही शाखेवरील
  पूर्वचिन्हांच्या उंचीची मर्यादा ठरते. प्रत्येक पूर्वचिन्हावर फक्त
  प्रारंभीच्या सूत्रांची सांत चिन्हांकित उपसूत्रे येऊ शकतात.
  प्रत्येक नवेपणा मागणाऱ्या चिन्हांकित मोडल सूत्रासाठी एकच
  साक्षीदार पूर्वचिन्ह पुरते, म्हणून प्रत्येक पूर्वचिन्हाची संततीही
  सांत असते. मर्यादित उंची आणि प्रत्येक ठिकाणी सांत संतती
  यांमुळे शाखेवरील सर्व पूर्वचिन्हे सांत संख्येची असतात.
  उर्वरित प्रत्येक नियमाचे आवश्यक अनुप्रयोगही सांत असतात.
  रचनेमुळे उरलेली प्रत्येक खुली शाखा संपूर्ण असते.
\end{proof}

\begin{thm}[संपूर्णता]
  \ollabel{thm:tableau-completeness}
  $\Gamma$ हा $1$ या एकाच पूर्वचिन्हाने चिन्हांकित सूत्रांचा संच असू द्या.
  $\Gamma$ साठी बंद !!{tableau} नसेल,
  तर $\Gamma$ पूर्ततायोग्य आहे.
\end{thm}

\begin{proof}
$\Gamma$ रिक्त असेल, तर एक जग असलेले कोणतेही प्रतिमान
ते पूर्ण करते. प्रथम $\Gamma$ रिक्त नसलेला सांत संच आहे
असे समजू. वरील विधानानुसार, $\Gamma$ साठी असा
!!a{tableau} असतो की त्याची प्रत्येक शाखा
बंद किंवा संपूर्ण असते. बंद !!{tableau} नसल्याने,
त्या !!a{tableau} मध्ये किमान एक शाखा
खुली आणि संपूर्ण असते. त्या शाखेवरील
पूर्वचिन्हयुक्त !!{formula}s चा संच
$\Delta$ आणि त्यात आलेल्या पूर्वचिन्हांचा
संच $P(\Delta)$ असू द्या.

आता $\mModel{M(\Delta)} = \tuple{P(\Delta), R, V}$
हे प्रतिमान ठरवू. त्यातील जगे म्हणजे~$\Delta$
मधील पूर्वचिन्हे, आणि प्राप्यता संबंध असा आहे:
\[
R\sigma\sigma' \quad \text{तेव्हा आणि तेव्हाच} \quad
\sigma'=\sigma.n \quad \text{काही~$n$ साठी}
\]
आणि
\[
V(p) = \Setabs{\sigma}{\sFmla{\True}{p}[\sigma] \in \Delta}.
\]
आता~$!A$ च्या रचनेवरून आगमनाने पुढील
दाखवू: $\sFmla{\True}{!A}[\sigma] \in
\Delta$ असेल, तर
$\mSat{M(\Delta)}{!A}[\sigma]$; आणि
$\sFmla{\False}{!A}[\sigma] \in \Delta$
असेल, तर $\mSat/{M(\Delta)}{!A}[\sigma]$.

\begin{enumerate}
  \item \indcase{!A}{p}{$\sFmla{\True}{\indfrm}[\sigma] \in \Delta$
    असेल, तर~$V$ च्या व्याख्येनुसार
    $\sigma \in V(p)$; म्हणून
    $\mSat{M(\Delta)}{\indfrm}[\sigma]$.

    $\sFmla{\False}{\indfrm}[\sigma] \in \Delta$ असेल, तर
    $\sFmla{\True}{\indfrm}[\sigma] \notin \Delta$,
    कारण अन्यथा शाखा बंद झाली असती.
    म्हणून $\sigma \notin V(p)$, आणि
    $\mSat/{M(\Delta)}{\indfrm}[\sigma]$.}
  \item \indcase{!A}{\lnot !B}{\iftag{probNot}{सरावासाठी.}{जर
      $\sFmla{\True}{\indfrm}[\sigma] \in \Delta$,
      तर शाखा संपूर्ण असल्याने
      $\sFmla{\False}{!B}[\sigma] \in \Delta$.
      आगमनाच्या गृहीतकानुसार
      $\mSat/{M(\Delta)}{!B}[\sigma]$;
      म्हणून $\mSat{M(\Delta)}{\indfrm}[\sigma]$.

      $\sFmla{\False}{\indfrm}[\sigma] \in \Delta$
      असेल, तर शाखा संपूर्ण असल्याने
      $\sFmla{\True}{!B}[\sigma] \in \Delta$.
      आगमनाच्या गृहीतकानुसार
      $\mSat{M(\Delta)}{!B}[\sigma]$;
      म्हणून $\mSat/{M(\Delta)}{\indfrm}[\sigma]$.}}
  \item \indcase{!A}{!B \land !C}{\iftag{probAnd}{सरावासाठी.}{जर
      $\sFmla{\True}{\indfrm}[\sigma] \in \Delta$,
      तर शाखा संपूर्ण असल्याने
      $\sFmla{\True}{!B}[\sigma] \in \Delta$
      आणि $\sFmla{\True}{!C}[\sigma] \in \Delta$
      ही दोन्ही. आगमनाच्या गृहीतकानुसार
      $\mSat{M(\Delta)}{!B}[\sigma]$ आणि
      $\mSat{M(\Delta)}{!C}[\sigma]$.
      म्हणून $\mSat{M(\Delta)}{\indfrm}[\sigma]$.

      $\sFmla{\False}{\indfrm}[\sigma] \in \Delta$
      असेल, तर शाखा संपूर्ण असल्याने
      $\sFmla{\False}{!B}[\sigma] \in \Delta$
      किंवा $\sFmla{\False}{!C}[\sigma] \in \Delta$.
      आगमनाच्या गृहीतकानुसार
      $\mSat/{M(\Delta)}{!B}[\sigma]$ किंवा
      $\mSat/{M(\Delta)}{!C}[\sigma]$.
      म्हणून $\mSat/{M(\Delta)}{\indfrm}[\sigma]$.}}
  \item \indcase{!A}{!B \lor !C}{\iftag{probOr}{सरावासाठी.}{जर
      $\sFmla{\True}{\indfrm}[\sigma] \in \Delta$,
      तर शाखा संपूर्ण असल्याने
      $\sFmla{\True}{!B}[\sigma] \in \Delta$
      किंवा $\sFmla{\True}{!C}[\sigma] \in \Delta$.
      आगमनाच्या गृहीतकानुसार
      $\mSat{M(\Delta)}{!B}[\sigma]$ किंवा
      $\mSat{M(\Delta)}{!C}[\sigma]$.
      म्हणून $\mSat{M(\Delta)}{\indfrm}[\sigma]$.

      $\sFmla{\False}{\indfrm}[\sigma] \in \Delta$
      असेल, तर शाखा संपूर्ण असल्याने
      $\sFmla{\False}{!B}[\sigma] \in \Delta$
      आणि $\sFmla{\False}{!C}[\sigma] \in \Delta$
      ही दोन्ही. आगमनाच्या गृहीतकानुसार
      $\mSat/{M(\Delta)}{!B}[\sigma]$ आणि
      $\mSat/{M(\Delta)}{!C}[\sigma]$.
      म्हणून $\mSat/{M(\Delta)}{\indfrm}[\sigma]$.}}
  \item \indcase{!A}{!B \lif !C}{\iftag{probIf}{सरावासाठी.}{जर
      $\sFmla{\True}{\indfrm}[\sigma] \in \Delta$,
      तर शाखा संपूर्ण असल्याने
      $\sFmla{\False}{!B}[\sigma] \in \Delta$
      किंवा $\sFmla{\True}{!C}[\sigma] \in \Delta$.
      आगमनाच्या गृहीतकानुसार
      $\mSat/{M(\Delta)}{!B}[\sigma]$ किंवा
      $\mSat{M(\Delta)}{!C}[\sigma]$.
      म्हणून $\mSat{M(\Delta)}{\indfrm}[\sigma]$.

      $\sFmla{\False}{\indfrm}[\sigma] \in \Delta$
      असेल, तर शाखा संपूर्ण असल्याने
      $\sFmla{\True}{!B}[\sigma] \in \Delta$
      आणि $\sFmla{\False}{!C}[\sigma] \in \Delta$
      ही दोन्ही. आगमनाच्या गृहीतकानुसार
      $\mSat{M(\Delta)}{!B}[\sigma]$ आणि
      $\mSat/{M(\Delta)}{!C}[\sigma]$.
      म्हणून $\mSat/{M(\Delta)}{\indfrm}[\sigma]$.}}
  \item \indcase{!A}{\Box !B}{\iftag{probBox}{सरावासाठी.}{जर
      $\sFmla{\True}{\indfrm}[\sigma] \in \Delta$,
      तर शाखा संपूर्ण असल्याने तिच्यावर
      वापरलेल्या प्रत्येक $\sigma.n$ साठी
      $\sFmla{\True}{!B}[\sigma.n] \in \Delta$.
      म्हणजे $R\sigma\sigma'$ होईल अशा
      प्रत्येक $\sigma' \in P(\Delta)$ साठी
      हे खरे आहे. आगमनाच्या गृहीतकानुसार,
      $R\sigma\sigma'$ असलेल्या प्रत्येक
      $\sigma'$ साठी
      $\mSat{M(\Delta)}{!B}[\sigma']$.
      म्हणून $\mSat{M(\Delta)}{\indfrm}[\sigma]$.

      $\sFmla{\False}{\indfrm}[\sigma] \in \Delta$
      असेल, तर शाखा संपूर्ण असल्याने
      काही $\sigma.n$ साठी
      $\sFmla{\False}{!B}[\sigma.n] \in \Delta$.
      आगमनाच्या गृहीतकानुसार
      $\mSat/{M(\Delta)}{!B}[\sigma.n]$.
      $R\sigma(\sigma.n)$ असल्याने
      $\mSat/{M(\Delta)}{!B}[\sigma']$
      होईल असे~$\sigma'$ असते. म्हणून
      $\mSat/{M(\Delta)}{\indfrm}[\sigma]$.}}
  \item \indcase{!A}{\Diamond !B}{\iftag{probDiamond}{सरावासाठी.}{जर
      $\sFmla{\True}{\indfrm}[\sigma] \in \Delta$,
      तर शाखा संपूर्ण असल्याने
      काही $\sigma.n$ साठी
      $\sFmla{\True}{!B}[\sigma.n] \in \Delta$.
      आगमनाच्या गृहीतकानुसार
      $\mSat{M(\Delta)}{!B}[\sigma.n]$.
      $R\sigma(\sigma.n)$ असल्याने
      $\mSat{M(\Delta)}{!B}[\sigma']$
      होईल असे~$\sigma'$ असते. म्हणून
      $\mSat{M(\Delta)}{\indfrm}[\sigma]$.

      $\sFmla{\False}{\indfrm}[\sigma] \in \Delta$
      असेल, तर शाखा संपूर्ण असल्याने तिच्यावर
      वापरलेल्या प्रत्येक $\sigma.n$ साठी
      $\sFmla{\False}{!B}[\sigma.n] \in \Delta$.
      म्हणजे $R\sigma\sigma'$ होईल अशा
      प्रत्येक $\sigma' \in P(\Delta)$ साठी
      हे खरे आहे. आगमनाच्या गृहीतकानुसार,
      $R\sigma\sigma'$ असलेल्या प्रत्येक
      $\sigma'$ साठी
      $\mSat/{M(\Delta)}{!B}[\sigma']$.
      म्हणून $\mSat/{M(\Delta)}{\indfrm}[\sigma]$.}}
\end{enumerate}
$\Gamma \subseteq \Delta$ असल्याने, प्रत्येक पूर्वचिन्हाला तेच जग
नेमणारे अर्थनिर्धारण घेऊन $\mSat{M(\Delta)}{\Gamma}[\mathrm{id}]$.
याने सांत प्रकरण सिद्ध झाले.
असांत $\Gamma$ साठी त्याच्या प्रत्येक सांत उपसंचाला बंद टॅब्लो
नसतो: तसा टॅब्लो असता, तर त्याची सांत गृहीतके $\Gamma$
मध्येही असल्याने तो $\Gamma$ साठी बंद टॅब्लो ठरला असता.
म्हणून सांत प्रकरणावरून प्रत्येक सांत उपसंच पूर्ततायोग्य आहे.
चिन्हांकित गृहीतकांचा साध्या मोडल सूत्रांचा संच असा ठरवू:
\[
\Upsilon = \Setabs{!B}{\sFmla{\True}{!B}[1]\in\Gamma}
\cup \Setabs{\lnot!B}{\sFmla{\False}{!B}[1]\in\Gamma}.
\]
याच्या प्रत्येक सांत उपसंचाला मूळ $\Gamma$ चा सांत उपसंच
पूर्ण करणाऱ्या प्रतिमानात एकाच जगावर सत्य करता येते.
त्यामुळे $\Upsilon$ हा \Log{K}-सुसंगत संच आहे: त्यावरून
असत्य सूत्र निष्पन्न होत असते, तर त्या सांत सिद्धतेतील सांत
गृहीतकांच्या पूर्ततायोग्यतेला \Log{K} ची निर्दोषता विरोधी ठरली
असती. \olref[nml][com][lin]{thm:lindenbaum} नुसार
$\Upsilon$ ला संपूर्ण \Log{K}-सुसंगत संच $\Lambda$ पर्यंत
वाढवता येते. \olref[nml][com][tru]{prop:truthlemma} नुसार
कॅनॉनिकल प्रतिमानात $\Lambda$ या जगावर $\Upsilon$ मधील
प्रत्येक सूत्र सत्य असते. $f(1)=\Lambda$ हे अर्थनिर्धारण मूळ
चिन्हांकित संच $\Gamma$ पूर्ण करते. याने असांत प्रकरणही सिद्ध झाले.
\end{proof}

\begin{prob}
\olref[nml][tab][cpl]{thm:tableau-completeness} चा
पुरावा पूर्ण करा.
\end{prob}

\begin{cor}
\ollabel{cor:entailment-completeness}
$\Gamma \Entails !A$ असेल, तर
$\Gamma \Proves !A$.
\end{cor}

\begin{cor}
\ollabel{cor:weak-completeness}
$!A$ सर्व प्रतिमानांमध्ये सत्य असेल,
तर $\Proves !A$.
\end{cor}

\end{document}
